%! AP Statistics — Unit 2, Lesson 2.5 — Student Packet Cover Sheet
\documentclass[10pt]{article}
\usepackage{apstats-article}
\usepackage{apstats-boxes}
\usepackage{ltablex}
\keepXColumns

\begin{document}

% Full-bleed navy banner
\begin{tikzpicture}[remember picture, overlay]
  \fill[navy] (current page.north west)
    rectangle ([yshift=-0.9in]current page.north east);
\end{tikzpicture}
\vspace{-0.8in}

\begin{center}
  {\color{white}\textbf{\LARGE AP Statistics}}\\[4pt]
  {\color{skymid}\large\textbf{Unit 2: Exploring Two-Variable Data}}\\[6pt]
  {\color{white}\textbf{\large Lesson 2.5 \quad Correlation}}
\end{center}

\vspace{0.22in}
\namedateperiod
\vspace{0.18in}

\begin{learningtargetbox}
\small
\begin{enumerate}[leftmargin=1.5em, itemsep=5pt]
  \item \ldots{} calculate the correlation coefficient $r$ for a bivariate dataset using
        technology. \hfill\textcolor{slate}{\small(DAT-1.B)}
  \item \ldots{} interpret $r$: state the direction, strength, and linear nature of the
        association in context. \hfill\textcolor{slate}{\small(DAT-1.C)}
  \item \ldots{} explain that $r$ is unit-free and always between $-1$ and $1$, inclusive.
        \hfill\textcolor{slate}{\small(DAT-1.C.1)}
  \item \ldots{} explain why correlation does not imply causation.
        \hfill\textcolor{slate}{\small(DAT-1.C.2)}
\end{enumerate}
\end{learningtargetbox}

\vspace{0.14in}

\begin{tocbox}
\small
\renewcommand{\arraystretch}{1.5}
\begin{tabularx}{\linewidth}{@{} c l X r @{}}
\rowcolor{sky}
\textbf{\#} & \textbf{Component} & \textbf{Description} & \textbf{Score} \\
\midrule
1 & Warm-Up        & Spiral review --- DOFS and estimating $r$ from scatterplots      & \blank{1.2cm} \\
2 & Guided Notes   & The correlation coefficient: formula, properties, interpretation  & \blank{1.2cm} \\
3 & Group Activity & Calculate and interpret $r$ for real-world datasets               & \blank{1.2cm} \\
4 & Exit Ticket    & Interpret a given $r$ and address causation                       & \blank{1.2cm} \\
5 & Homework       & Multi-context correlation practice and scatterplot matching        & \blank{1.2cm} \\
\midrule
  & \multicolumn{2}{r}{\textbf{Total}} & \blank{1.2cm} \\
\end{tabularx}
\end{tocbox}

\vspace{0.14in}

\begin{remindbox}
\small
The correlation coefficient $r$ measures the \emph{direction} and \emph{strength} of a
\textbf{linear} association. Keep these ideas in mind.

\vspace{6pt}
\begin{tabularx}{\linewidth}{@{} >{\centering\arraybackslash\columncolor{goldbg}}p{0.06\linewidth}
                                    X @{}}
\textbf{\large!} &
\textbf{$r$ only measures \emph{linear} association.}\enspace
A value of $r \approx 0$ means no linear association ---
a strong curved relationship could still exist. Always look at the scatterplot. \\[6pt]
\textbf{\large!} &
\textbf{$r$ is unit-free and symmetric.}\enspace
Switching $x$ and $y$ does not change $r$. Multiplying a variable by a positive
constant does not change $r$. \\[6pt]
\textbf{\large!} &
\textbf{Interpret in context.}\enspace
Always name the variables: ``There is a strong negative linear association between
engine size and fuel efficiency'' --- not just ``$r = -0.87$.'' \\[6pt]
\textbf{\large!} &
\textbf{Correlation $\neq$ causation.}\enspace
A high $|r|$ does \emph{not} mean one variable causes the other.
A lurking variable or coincidence may explain the association.
\end{tabularx}
\end{remindbox}

\end{document}
